\begin{recipe}{Arretjescake}
\kicker[Dessert,Zoet,Nederlands]{Dessert · zoet · Nederlands}
\index[register]{Arretjescake}
\index[register]{Chocolade}
\index[register]{Cacao}
\index[register]{Boter}

\meta{15 minuten + minstens 2 uur koelen \dvd Voor 10 personen \dvd Eenvoudig}

\lettrine{A}{rretjescake} is een ouderwetse chocoladecake zonder ei. Je maakt
hem van gebroken koekjes in een smeuïge chocolademassa, zonder dat er een oven
aan te pas komt. Smelt de chocolade au bain-marie en laat de cake daarna
opstijven in de koelkast.

\blockrule{Ingrediënten}
\begin{ingredients}
  \ing{250 g}{boter of margarine}\index[register]{Boter}
  \ing{100 g}{pure chocolade}\index[register]{Chocolade}
  \ing{50 g}{cacao}\index[register]{Cacao}
  \ing{200 g}{basterdsuiker}
  \ing{1 zakje}{vanillesuiker (of een scheutje vanille-extract)}
  \ing{200 g}{koekjes (bijvoorbeeld Marie)}
\end{ingredients}

\blockrule{Bereiding}
\tip{Een beetje oploskoffie door het chocolademengsel geeft de cake een sterkere
smaak. Roomboterspritsen in plaats van (of naast) de koekjes maken hem steviger
van bite.}
\begin{steps}
  \item Bekleed een cakeblik met 2 lagen huishoudfolie, zodat je de cake er straks
        makkelijk uit kunt halen.
  \item Breek de koekjes in kleine stukjes en de chocolade in kleinere stukken.
        Snijd de boter in blokjes.
  \item Zet een pan met een bodempje water op het vuur en breng aan de kook. Doe
        ondertussen de boter, chocolade, cacao, basterdsuiker en vanillesuiker samen
        in een hittebestendige kom.
  \item Zodra het water kookt, zet je het vuur laag en zet je de kom op de pan
        (au bain-marie), zonder dat de bodem het water raakt. Roer op laag vuur tot
        de boter en chocolade volledig gesmolten zijn en het mengsel glad en
        glanzend is, dit duurt ongeveer 10 minuten.
  \item Haal de kom van het vuur en schep de gebroken koekjes erdoor. Meng goed tot
        de koekjes gelijkmatig verdeeld zijn door het chocolademengsel.
  \item Schep het mengsel in het beklede cakeblik en druk het geheel goed aan. Vouw
        de huishoudfolie over de bovenkant en laat de cake minstens 2 uur in de
        koelkast opstijven.
  \item Haal de cake met behulp van de folie uit de vorm en snijd er plakken of
        punten van. Serveer de cake niet te koud en houd de porties klein, want hij
        is machtig.
\end{steps}
\end{recipe}
